\section{Computational certificates and review}\label{sec:computations}
The proof uses finite calculations at three points: the arithmetic
integral, the maximum of the real phase, and the nonzero seed.
Each is performed with exact integer or rational arithmetic. Decimal
values in the text are illustrations of the certified bounds.

\begin{center}
\begin{tabular}{@{}lll@{}}
\toprule
Calculation & Size & Certified conclusion \\
\midrule
Arithmetic partition & $38\,371$ intervals & $A_{\rm cert}<41/5$ \\
Critical polynomial & degree $33$ & $\max\Phi<-18.83$ \\
Seed coefficients & $2077$ at each of two primes & $\widetilde q\ne0$ \\
\bottomrule
\end{tabular}
\end{center}

The program \texttt{certify\_fifteen.py} generates the rational event
partition and computes its maxima by sweeping the jumps in $y$.
A separate C++ program, \texttt{verify\_fifteen\_arithmetic.cpp},
recomputes every supplied midpoint maximum by direct floor evaluation
on the cells. Thus the second calculation checks the maxima; completeness
of the event partition follows from the affine-offset argument in
Section~\ref{sec:decay} and its explicit enumeration in the first program.

The polynomial calculation in \texttt{certify\_phase\_polynomial.py}
multiplies the factors in \eqref{eq:critical-polynomial}, checks every
coefficient sign, brackets its positive root, and encloses the logarithms
used in \eqref{eq:phase-certificate}. The arithmetic program also
checks the weaker phase bound $-18.83$ by a separate concavity argument.

For the seed, \texttt{seed\_fifteen.py} uses truncated products over
finite fields. The separate calculation
\texttt{verify\_fifteen\_seed.py} uses the Taylor recurrence described
after \eqref{eq:seed-residues}. It computes all coefficients before
comparing them with the saved product calculation. Both give the
residues in \eqref{eq:seed-residues}. The scripts, certificate data,
and reproduction commands accompany the source in \texttt{checks/}.

The work was prepared with OpenAI Codex using GPT-6-Astra max.
The isolated-pole refinement and the fifteen-value conclusion were
obtained separately in two research sessions using the same public
sources, without access to each other's unpublished work. This paper
combines the subsequent comparison into one proof. The supplied
calculations have been replayed for this manuscript.

An independent mathematical audit of the complete fourteen-page version,
reported on 11 November 2026, found no mathematical error or gap affecting
Theorem~\ref{thm:main}. It checked the general arguments and independently
reproduced the arithmetic integral, the phase bound, and all $2077$ seed
coefficients at both primes. The audit accepts the stated classical inputs,
including Stirling's formula and the prime number theorem; it is not
formal verification. The accompanying audit record identifies the
reviewed version and the verification scope.
